%-------------------------
% Resume in Latex
% Author : Jake Gutierrez
% Based off of: https://github.com/sb2nov/resume
% License : MIT
%------------------------

\documentclass[letterpaper,10pt]{article}

\usepackage{iftex}
\ifPDFTeX
  \usepackage[T2A,T1]{fontenc}
  \usepackage[utf8]{inputenc}
  \usepackage[english,russian]{babel}
  \input{glyphtounicode}
  \pdfgentounicode=1
\else
  \usepackage{fontspec}
  \usepackage{polyglossia}
  \setmainlanguage{english}
  \setotherlanguage{russian}
  \setmainfont{Times New Roman}
\fi

\usepackage{latexsym}
\usepackage[empty]{fullpage}
\usepackage{titlesec}
\usepackage{marvosym}
\usepackage[usenames,dvipsnames]{color}
\usepackage{verbatim}
\usepackage{enumitem}
\usepackage[hidelinks]{hyperref}
\usepackage{fancyhdr}
\usepackage{tabularx}

\pagestyle{fancy}
\fancyhf{}
\fancyfoot{}
\renewcommand{\headrulewidth}{0pt}
\renewcommand{\footrulewidth}{0pt}

\addtolength{\oddsidemargin}{-0.5in}
\addtolength{\evensidemargin}{-0.5in}
\addtolength{\textwidth}{1in}
\addtolength{\topmargin}{-.62in}
\addtolength{\textheight}{1.24in}

\urlstyle{same}

\raggedbottom
\raggedright
\setlength{\tabcolsep}{0in}

\titleformat{\section}{
  \vspace{-4pt}\scshape\raggedright\large
}{}{0em}{}[\color{black}\titlerule \vspace{-5pt}]

\newcommand{\resumeItem}[1]{
  \item\small{
    {#1 \vspace{-2pt}}
  }
}

\newcommand{\resumeSubheading}[4]{
  \vspace{-1pt}\item
    \begin{tabular*}{0.97\textwidth}[t]{l@{\extracolsep{\fill}}r}
      \textbf{#1} & #2 \\
      \textit{\small#3} & \textit{\small #4} \\
    \end{tabular*}\vspace{-7pt}
}

\newcommand{\resumeProjectHeading}[2]{
    \item
    \begin{tabular*}{0.97\textwidth}{l@{\extracolsep{\fill}}r}
      \small#1 & #2 \\
    \end{tabular*}\vspace{-7pt}
}

\newcommand{\resumeSubItem}[1]{\resumeItem{#1}\vspace{-4pt}}

\renewcommand\labelitemii{$\vcenter{\hbox{\tiny$\bullet$}}$}

\newcommand{\resumeSubHeadingListStart}{\begin{itemize}[leftmargin=0.15in, label={}]}
\newcommand{\resumeSubHeadingListEnd}{\end{itemize}}
\newcommand{\resumeItemListStart}{\begin{itemize}[itemsep=0pt, parsep=0pt, topsep=1pt]}
\newcommand{\resumeItemListEnd}{\end{itemize}\vspace{-5pt}}

\begin{document}

\begin{center}
    \textbf{\Huge \scshape Andrey Urakov} \\ \vspace{1pt}
    \small{}Yekaterinburg $|$ \href{https://t.me/andrey\_urakov\_ml}{\underline{tg: @andrey\_urakov\_ml}} $|$ \href{mailto:urakov18.a@gmail.com}{\underline{urakov18.a@gmail.com}} $|$ \href{https://github.com/Andrey-Good}{\underline{github.com/Andrey-Good}}\end{center}

\section{Education}
  \resumeSubHeadingListStart
    \resumeSubheading
      {UrFU (Ural Federal University)}{Yekaterinburg}
      {Bachelor's degree, Applied Artificial Intelligence (2nd year)}{2024 -- 2028}
  \resumeSubHeadingListEnd

\section{Relevant Coursework}
 \begin{itemize}[leftmargin=0.15in, label={}]
    \small{\item{Nand2Tetris, Advanced Programming, Algorithms and Complexity Analysis, Discrete Mathematics and Mathematical Logic, Probability and Mathematical Statistics, Mathematical Analysis, Vector Analysis, Machine Learning, Databases, Non-relational Databases, Computer Architecture, Computer Networks, Data Analysis \& AI}}
 \end{itemize}

\section{Experience}
  \resumeSubHeadingListStart
    \resumeSubheading
      {IDScan}{Chelyabinsk}
      {Junior Data Annotator (Full-time / Summer contract)}{July 2022 -- August 2022}
      \resumeItemListStart
        \resumeItem{Worked offline in a closed environment with sensitive personal data that the ML team later used to train CV models.}
        \resumeItem{Performed manual annotation and transcription of images to create an OCR dataset.}
        \resumeItem{Performed data cleaning and noise filtering according to technical quality criteria.}
      \resumeItemListEnd
  \resumeSubHeadingListEnd

\section{Projects (R\&D / Engineering)}
    \resumeSubHeadingListStart
      \resumeProjectHeading
          {\textbf{Active Learning SDK} $|$ \emph{Product Logic, Active Learning, Evaluation, Lead (Team Project)} $|$ \href{https://github.com/Andrey-Good/active-learning-orchestrator}{\underline{[Source Code]}}}{2026}
          \resumeItemListStart
            \resumeItem{Used AI-agent-assisted development while owning the project idea, product contract, SDK workflow, evaluation design, and acceptance criteria.}
            \resumeItem{Selected acquisition strategies for the SDK and benchmark suite, balancing uncertainty, diversity, class/group balance, and random baselines.}
            \resumeItem{Researched calibration and Temperature Scaling; translated useful findings into benchmark criteria with ECE/NLL/Brier metrics and calibration-aware stop checks.}
            \resumeItem{Led evaluation and evidence work on capped real text-classification datasets: quality gate PASS; Banking77 adaptive acquisition matched capped full-train macro-F1 at 400 labels.}
          \resumeItemListEnd
      \resumeProjectHeading
          {\textbf{Local-First AI Agent Hub} $|$ \emph{Architecture Author, SQLite, Lead (Team of 5)} $|$ \href{https://drive.google.com/drive/folders/13iSz2zMOBYBxjSZVT6gQPQSsaGLZ3evt?usp=sharing}{\underline{[Arch Showcase]}}}{2025}
          \resumeItemListStart
            \resumeItem{Defined functional scope, roadmap, and task decomposition; owned delivery timing and quality.}
            \resumeItem{Authored the Blueprint + Spec: defined the architecture, business logic, and reliability acceptance criteria.}
            \resumeItem{Designed a failure-recovery mechanism: the system restores state consistency after a sudden process crash without data loss.}
            \resumeItem{Designed a task scheduler with protection against parallel-run conflicts and stuck tasks, using atomic operations and SQLite WAL.}
            \resumeItem{Designed an isolated execution environment based on uv to prevent dependency conflicts.}
            \resumeItem{Documented the Reconcile Loop pattern for trigger management in the specification: automatically bringing the system to the target state.}
            \resumeItem{Designed an event-driven architecture: separated business logic from the API.}
            \resumeItem{Designed I/O performance optimization through hybrid storage: metadata is stored in the database, while heavy artifacts are stored in the file system.}
            \resumeItem{UX: designed the UX architecture, user flow, and states; without a full UI design.}
          \resumeItemListEnd
      \resumeProjectHeading
          {\textbf{Melanoma Detection System} $|$ \emph{ML Engineer, PyTorch, ONNX, Lead (Team of 4)}}{2025}
          \resumeItemListStart
            \resumeItem{Coordinated the full development cycle: defined functional scope and roadmap, decomposed requirements into technical tasks with acceptance criteria, and ensured delivery timing and quality. Organized the project defense (UrFU project work): 1st place among teams from two cohorts, 98/100.}
            \resumeItem{Implemented the full training cycle for CNN-based image classification.}
            \resumeItem{Applied data-quality improvement techniques: augmentation and class-imbalance mitigation.}
            \resumeItem{Optimized the model for mobile inference: ONNX conversion and quantization.}
            \resumeItem{UI/UX: authored the concept and UI design; designed interface behavior (states/errors) and user flow.}
          \resumeItemListEnd
      \resumeProjectHeading
          {\textbf{Review Summarizer Extension} $|$ \emph{Backend \& Core Logic, NLP, JS, NodeJS, Lead (Team of 4)} $|$ \href{https://github.com/Andrey-Good/RevAI}{\underline{[Source Code]}}}{2024}
          \resumeItemListStart
            \resumeItem{Coordinated the full development cycle: defined functional scope and roadmap, decomposed requirements into technical tasks with acceptance criteria, and ensured delivery timing and quality.}
            \resumeItem{Developed the full extension architecture (Manifest v3): core structure, a custom DOM parser, scripts, and binding the logic to the HTML interface.}
            \resumeItem{Implemented the server side: proxying requests to an LLM API and caching through request hashing to reduce latency and costs.}
          \resumeItemListEnd
    \resumeSubHeadingListEnd

\section{Skills}
 \begin{itemize}[leftmargin=0.15in, label={}]
    \small{\item{
     \textbf{Languages}{: Python (Multiprocessing, AsyncIO, Typer), SQL (SQLite)} \\     \textbf{ML/Data}{: PyTorch, Classical ML (linear models, SVM/kernels, Bayesian methods, PCA), Model Evaluation, Data Annotation} \\     \textbf{Engineering}{: Git, System Design (Local-First systems, Event-driven), Architecture Patterns (IPC, Queues), Linux (WSL environment)} \\     \textbf{Soft Skills}{: Technical Writing (RFC/Specs), Project Leadership (5 people), Research Mindset} \\     \textbf{Human Languages}{: Russian (Native), English (B1 - Intermediate)}    }}
 \end{itemize}

\end{document}
